\documentclass[12pt,a4paper]{article}

% Required Packages
\usepackage[utf8]{inputenc}
\usepackage{geometry}
\geometry{a4paper, margin=1in}
\usepackage{graphicx}
\usepackage{hyperref}
\usepackage{xcolor}
\usepackage{titlesec}
\usepackage{enumitem}

% Setup Hyperlinks
\hypersetup{
    colorlinks=true,
    linkcolor=blue,
    filecolor=magenta,      
    urlcolor=cyan,
    pdftitle={Real-Time Stocks Dashboard: Technical Architecture Report},
}

% Title formats
\titleformat{\section}{\Large\bfseries\color{darkgray}}{\thesection}{1em}{}
\titleformat{\subsection}{\large\bfseries\color{gray}}{\thesubsection}{1em}{}

% Custom Colors
\definecolor{darkgray}{RGB}{50,50,50}
\definecolor{gray}{RGB}{100,100,100}

\begin{document}

\title{\textbf{Real-Time Stocks: Pipeline Evolution (V1 to V2)} \\ \Large Technical Architecture \& Implementation Report}
\author{Samarth Agarwal \\ \href{mailto:samarthagrawal526@gmail.com}{samarthagrawal526@gmail.com}}
\date{\today}

\maketitle

\begin{abstract}
This report outlines the technical evolution of a real-time financial market data system. Originally conceived as an enterprise-grade Modern Data Stack (Version 1) utilizing Kafka, Snowflake, and Airflow, the project has transitioned into an ultra-low-latency, infinitely scalable Global Serverless Edge application (Version 2) utilizing the Cloudflare ecosystem and React. This document details the architectural shift, the strategic on-demand caching model of Version 2, and the resulting performance improvements.
\end{abstract}

\vspace{1cm}
\tableofcontents
\newpage

\section{Evolution Overview}
The project transitioned through two distinct architectural phases:
\begin{itemize}
    \item \textbf{Version 1 (Legacy):} A heavyweight Modern Data Stack (MDS) utilizing containerized orchestration (Docker), stream processing (Kafka), and cloud warehousing (Snowflake).
    \item \textbf{Version 2 (Current):} A serverless, edge-computed dashboard hosted on Cloudflare Pages/Workers, optimizing for cost, latency, and global availability. \textit{Live Demo: \url{https://d0fc8457.marketpulse-91t.pages.dev}}
\end{itemize}

\section{Version 1: The Heavy Modern Data Stack}
Version 1 focused on absolute data persistence and enterprise-grade ETL pipelines. It utilized Apache Kafka to stream live symbols into a MinIO-based data lake, followed by Apache Airflow orchestrating dbt transformations inside Snowflake. While statistically robust, Version 1 required always-on compute resources and significant maintenance overhead.

\section{Version 2: Serverless Edge Architecture}
Version 2 represents the matured, cost-efficient state of the project. By shedding local orchestration, the application now runs entirely on the Cloudflare global network.

\subsection{The Frontend (React + Vite)}
The user-facing application was built using \textbf{React} and the \textbf{Vite} toolset.
\begin{itemize}
    \item \textbf{State Management:} Features real-time sync toggles and timeframe analysis (1H, 1D, 1W) with automatic data staleness detection.
    \item \textbf{Deployment:} Deployed via \textbf{Cloudflare Pages} for instant global delivery.
\end{itemize}

\subsection{The Backend / API (Cloudflare Workers)}
All backend logic is handled by a standalone \textbf{Cloudflare Worker} proxy. Operating on V8 Isolates, it provides zero-millisecond cold starts and proxies all API endpoints natively without the need for a dedicated server instance.

\subsection{Database \& Storage}
\begin{itemize}
    \item \textbf{Cloudflare D1 (SQLite):} Serves as the primary edge-database, logging fetched transactions locally to provide sub-10ms cache hits.
    \item \textbf{Cloudflare R2 (Object Storage):} Acts as the archival data lake for raw financial JSON payloads.
\end{itemize}

\section{Architectural Innovation: On-Demand Cache}
To prevent API exhaustion, the Version 2 Worker implements a smart-caching proxy. Data is only fetched from the external provider if the local D1 cache is older than 60 seconds. This ensures that even under heavy concurrent load (e.g., 10,000+ users), the upstream provider is only queried once per minute per symbol, while users receive instantaneous responses from the edge database.

\section{Conclusion}
The migration from Version 1's monolithic enterprise stack to Version 2's serverless edge architecture highlights a shift toward hyper-efficiency. Version 2 achieves superior user responsiveness and feature density while maintaining a zero-maintenance, cost-optimized operational profile.

\end{document}
